\documentclass[11pt]{article}
\usepackage[margin=1in]{geometry}
\usepackage{setspace}
\onehalfspacing

% ── Mathematics ──────────────────────────────────────────────
\usepackage{amsmath,amssymb,amsthm}
\usepackage{mathtools}

% ── Tables and figures ───────────────────────────────────────
\usepackage{booktabs}
\usepackage{array}
\usepackage{graphicx}
\usepackage{caption}
\captionsetup{font=small}

% ── Lists ────────────────────────────────────────────────────
\usepackage{enumitem}

% ── Colors and links ─────────────────────────────────────────
\usepackage[colorlinks=true,linkcolor=blue,citecolor=blue,urlcolor=blue]{hyperref}

% ── Bibliography ─────────────────────────────────────────────
\usepackage[numbers,sort&compress]{natbib}

% ── Theorem environments ─────────────────────────────────────
\newtheorem{theorem}{Theorem}[section]
\newtheorem{proposition}[theorem]{Proposition}
\newtheorem{lemma}[theorem]{Lemma}
\newtheorem{corollary}[theorem]{Corollary}
\newtheorem{definition}[theorem]{Definition}
\newtheorem{remark}[theorem]{Remark}

% ── Custom operators ─────────────────────────────────────────
\DeclareMathOperator{\TIP}{TIP}
\DeclareMathOperator{\supp}{supp}
\DeclareMathOperator{\Rips}{Rips}
\DeclareMathOperator{\diam}{diam}
\DeclareMathOperator{\Tr}{Tr}
\DeclareMathOperator*{\argmin}{arg\,min}

% ── Supplement numbering ─────────────────────────────────────
\renewcommand{\thesection}{S\arabic{section}}
\renewcommand{\thetable}{S\arabic{table}}
\renewcommand{\thefigure}{S\arabic{figure}}
\renewcommand{\theequation}{S\arabic{equation}}

\emergencystretch=1em

\title{Supplementary Information:\\
CC-SPR: Geometry-Sensitive Sparse Persistent Representatives\\
for Interpretable Omics Feature Selection}

% PLACEHOLDER: Match main manuscript authors
\author{Author One$^{1,\ast}$, Author Two$^{1}$, Author Three$^{1}$\\[4pt]
\small $^{1}$[Department], [Institution], [City], [Country]\\
\small $^\ast$Corresponding author: \href{mailto:author@institution.edu}{author@institution.edu}}

\date{}

\begin{document}
\maketitle

\tableofcontents
\newpage

% ============================================================
\section{Mathematical Background}
\label{sec:supp_background}
% ============================================================

This section provides the formal mathematical definitions underlying the CC-SPR framework.
These are condensed in the main text (Section~2) and presented here in full detail.

\subsection{Simplicial homology and persistence}

Let $X = \{x_1, \ldots, x_n\} \subset \mathbb{R}^p$ be a point cloud equipped with a metric $d$.
The \emph{Vietoris--Rips complex} at scale $\varepsilon$ is
\[
\Rips_\varepsilon(X, d) = \bigl\{ \sigma \subseteq X : \diam_d(\sigma) \le \varepsilon \bigr\}.
\]
As $\varepsilon$ increases from $0$ to $\infty$, the nested family $\{\Rips_\varepsilon\}_{\varepsilon \ge 0}$ forms a filtration.
Applying simplicial homology with field coefficients $\mathbb{F}$ to this filtration yields a persistence module.
By the structure theorem for graded modules over a PID, this module decomposes into a direct sum of interval modules, each corresponding to an interval $[b_i, d_i)$ in the barcode \cite{zomorodian2005}.

\begin{definition}[Persistence barcode]\label{def:supp_barcode}
The \emph{barcode} of a filtered simplicial complex is the multiset of intervals $\{[b_i, d_i)\}_{i \in I}$ encoding the birth and death scales of homology classes.
The \emph{persistence} (or \emph{lifetime}) of the $i$-th feature is $\pi_i = d_i - b_i$, and its \emph{midlife} is $m_i = (b_i + d_i)/2$.
\end{definition}

\begin{definition}[Chain groups and boundary maps]\label{def:supp_chains}
Let $K$ be a simplicial complex.
The \emph{$k$-chain group} $C_k(K; \mathbb{F})$ is the $\mathbb{F}$-vector space with basis the $k$-simplices of $K$.
The \emph{boundary operator} $\partial_k: C_k \to C_{k-1}$ is the linear map defined on each $k$-simplex $\sigma = [v_0, \ldots, v_k]$ by
\[
\partial_k(\sigma) = \sum_{i=0}^k (-1)^i [v_0, \ldots, \hat{v}_i, \ldots, v_k],
\]
where $\hat{v}_i$ denotes omission of vertex $v_i$.
The fundamental identity $\partial_{k-1} \circ \partial_k = 0$ ensures that $\operatorname{im}(\partial_{k+1}) \subseteq \ker(\partial_k)$, and the $k$-th homology group is $H_k(K; \mathbb{F}) = \ker(\partial_k) / \operatorname{im}(\partial_{k+1})$.
\end{definition}

\subsection{Harmonic representatives and the dense support problem}

For a selected $H_1$ bar $[b, d)$, let $z_0 \in C_1(\Rips_\varepsilon; \mathbb{R})$ be an initial cycle representative at scale $\varepsilon \in [b, d)$, and let $\partial_2 : C_2 \to C_1$ denote the boundary operator mapping 2-chains (triangles) to 1-chains (edges).
The full set of cycles homologous to $z_0$ forms the \emph{affine homology class}
\begin{equation}\label{eq:supp_affine}
[z_0] + \operatorname{im}(\partial_2) = \{z_0 + \partial_2 y : y \in C_2\}.
\end{equation}

The \emph{harmonic representative} minimizes the $\ell_2$ energy within this affine class:
\begin{equation}\label{eq:supp_harmonic}
\hat{z}_{\mathrm{harm}} = \argmin_{y \in C_2} \| z_0 + \partial_2 y \|_2^2.
\end{equation}
This is a standard least-squares problem whose solution is the orthogonal projection of $z_0$ onto $\ker(\partial_1) \cap (\operatorname{im} \partial_2)^\perp$, coinciding with the harmonic representative in the Hodge-theoretic sense when the chain complex is equipped with the standard inner product \cite{harmonicPH, friedman1998}.

\begin{remark}[Dense support problem]\label{rem:supp_dense}
The solution to \eqref{eq:supp_harmonic} typically has $\|\hat{z}_{\mathrm{harm}}\|_0 = O(|E|)$ where $|E|$ is the number of edges at scale $\varepsilon$.
In omics applications where $n$ ranges from $10^2$ to $10^4$, this can mean hundreds or thousands of nonzero edge weights, diluting the biological signal when projected to feature space.
The harmonic representative finds the smoothest cycle, not the most parsimonious one.
\end{remark}


% ============================================================
\section{Proofs of Convexity and Sparsity Guarantees}
\label{sec:supp_proofs}
% ============================================================

This section provides the formal proofs referenced in the main text (Section~2.2).

\subsection{Convexity of the CC-SPR objective}

\begin{proposition}[Convexity]\label{prop:supp_convex}
The CC-SPR objective
\begin{equation}\label{eq:supp_ccspr}
f(y) = \| z_0 + \partial_2 y \|_2^2 + \lambda_1 \| z_0 + \partial_2 y \|_1 + \lambda_2 \| y \|_2^2
\end{equation}
is convex in $y$. When $\lambda_2 > 0$, it is strongly convex with a unique minimizer.
\end{proposition}

\begin{proof}
Each term is a convex function of $y$:
\begin{enumerate}[leftmargin=1.5em]
\item $f_1(y) = \|z_0 + \partial_2 y\|_2^2$ is convex as a composition of a convex quadratic with the affine map $y \mapsto z_0 + \partial_2 y$;
\item $f_2(y) = \lambda_1 \|z_0 + \partial_2 y\|_1$ is convex because the $\ell_1$ norm is convex and composition with affine maps preserves convexity;
\item $f_3(y) = \lambda_2 \|y\|_2^2$ is convex (strongly convex when $\lambda_2 > 0$).
\end{enumerate}
A nonnegative sum of convex functions is convex, and strong convexity of any summand implies strong convexity of the sum.
\end{proof}

\begin{remark}[Connection to elastic net regression]\label{rem:supp_elasticnet}
The structure of \eqref{eq:supp_ccspr} mirrors the elastic net \cite{zou2005elasticnet}: $\lambda_1$ plays the role of the LASSO penalty promoting sparsity in the edge weights, while $\lambda_2$ provides the ridge-like stability that prevents ill-conditioning when $\partial_2$ is rank-deficient.
The key difference is that the ``design matrix'' here is the boundary operator $\partial_2$ rather than a data matrix, and the ``response'' is the initial cycle $z_0$.
\end{remark}

\subsection{Sparsity guarantee}

\begin{proposition}[Sparsity guarantee]\label{prop:supp_sparsity}
For any $\lambda_1 > 0$, the solution $\hat{z}_{\mathrm{spr}}$ satisfies
\[
\|\hat{z}_{\mathrm{spr}}\|_0 \le \|\hat{z}_{\mathrm{harm}}\|_0,
\]
with equality only when the harmonic representative already has minimal $\ell_1$ norm in its homology class.
\end{proposition}

\begin{proof}
The $\ell_1$ penalty acts as a soft-thresholding operator on the edge weights of the representative.
Any edge weight $|z_e| < \tau(\lambda_1)$ for an appropriate threshold $\tau$ is driven to zero.
Since the harmonic representative has no such thresholding, its support is at least as large.
Equality requires that every nonzero weight in $\hat{z}_{\mathrm{harm}}$ exceeds $\tau(\lambda_1)$, which occurs only when the harmonic solution already has minimal $\ell_1$ norm---a non-generic condition.
\end{proof}


% ============================================================
\section{Detailed Support Diagnostics Definitions}
\label{sec:supp_diagnostics}
% ============================================================

Beyond TIP (defined in the main text, Section~2.3), the quality of a sparse representative can be characterized by several support diagnostics used throughout the results.

\begin{definition}[Support size]\label{def:supp_support}
For a representative $\hat{z} \in \mathbb{R}^{|E|}$, the \emph{support size} is $\|\hat{z}\|_0 = |\{e \in E : \hat{z}_e \ne 0\}|$.
After projection to feature space, the \emph{projected support size} is $|\{j \in \{1,\ldots,p\} : |w_j| > 0\}|$ where $w \in \mathbb{R}^p$ is the projected feature-importance vector.
\end{definition}

\begin{definition}[Representative entropy]\label{def:supp_entropy}
Let $\hat{p}_e = |\hat{z}_e| / \sum_{e'} |\hat{z}_{e'}|$ be the normalized absolute edge weights.
The \emph{representative entropy} is
\[
H(\hat{z}) = -\sum_{e : \hat{z}_e \ne 0} \hat{p}_e \log \hat{p}_e.
\]
Low entropy indicates concentrated support (weight localized on a few edges); high entropy indicates diffuse support (weight spread across many edges).
\end{definition}

\begin{definition}[Gini coefficient of representative weights]\label{def:supp_gini}
The \emph{Gini coefficient} of the weight distribution provides a complementary inequality measure:
\[
G(\hat{z}) = \frac{\sum_{i=1}^m \sum_{j=1}^m ||\hat{z}_{e_i}| - |\hat{z}_{e_j}||}{2m \sum_{i=1}^m |\hat{z}_{e_i}|},
\]
where $m = \|\hat{z}\|_0$ is the support size.
High Gini indicates that weight is concentrated on a few edges; low Gini indicates uniform distribution.
\end{definition}

\begin{remark}[Interpreting TIP entropy and Gini in results tables]\label{rem:supp_concentration}
In the main text Tables~3--4 and Table~7 of the main text, the ``TIP Ent.''\ and ``TIP Gini'' columns report the entropy and Gini coefficient of the TIP profile itself (not of the raw edge weights).
In omics, a representative with high Gini coefficient and low entropy concentrates its topological signal on a small number of sample-pair relationships, which project to a compact gene set.
This concentration is what makes the representative biologically actionable: it identifies a gene module rather than a diffuse background signal.
The $\ell_1$ penalty in the CC-SPR objective directly promotes this concentration by driving small weights to zero.
\end{remark}


% ============================================================
\section{Stability Proposition Under Bounded Metric Perturbation}
\label{sec:supp_stability}
% ============================================================

This section formalizes the theoretical guarantee that backend choice can be studied as a controlled perturbation of topological output, referenced in the main text (Section~2.4).

\begin{proposition}[Filtration stability under backend perturbation]\label{prop:supp_stability}
Let $D^{(a)}$ and $D^{(b)}$ be distance matrices from two backends.
Define $\varepsilon_\infty = \|D^{(a)} - D^{(b)}\|_\infty$.
Then the Rips filtrations satisfy an $\varepsilon_\infty$-interleaving:
\[
\Rips_r(X, D^{(a)}) \subseteq \Rips_{r + \varepsilon_\infty}(X, D^{(b)}) \subseteq \Rips_{r + 2\varepsilon_\infty}(X, D^{(a)})
\]
for all $r \ge 0$.
Consequently, the bottleneck distance between the corresponding persistence diagrams satisfies
\[
d_B\bigl(\mathrm{Dgm}(D^{(a)}),\, \mathrm{Dgm}(D^{(b)})\bigr) \le \varepsilon_\infty.
\]
\end{proposition}

\begin{proof}
The first inclusion follows because if $\diam_{D^{(a)}}(\sigma) \le r$, then $\diam_{D^{(b)}}(\sigma) \le r + \varepsilon_\infty$ by the definition of $\varepsilon_\infty = \max_{i,j} |D^{(a)}_{ij} - D^{(b)}_{ij}|$.
The second inclusion follows symmetrically.
The bottleneck bound is the standard algebraic stability theorem for Rips filtrations applied to $\varepsilon_\infty$-interleaved persistence modules \cite{chazal2014persistence}.
\end{proof}

\begin{remark}[Practical implication]\label{rem:supp_perturbation}
This proposition means that backend choice can be studied as a \emph{controlled perturbation} of the topological output.
When $\varepsilon_\infty$ is small relative to persistence values, the barcode is approximately preserved; when it is large, the backend is genuinely changing the topological signal.
The benchmark results in the main text can be understood through this lens: backends that produce substantially different distance matrices yield different topological features and different downstream classifier performance.
\end{remark}

\begin{remark}[Bridge edges, bottleneck sensitivity, and backend sharpening]\label{rem:supp_bridge}
Backends that preferentially upweight bottleneck or bridge-like edges---as Ricci flow does for negatively curved edges---can sharpen the persistent signal by increasing the filtration parameter at which spurious shortcuts create topological noise.
This is the mechanism by which curvature-aware backends can improve representative quality: by delaying the merger of distinct topological components, they preserve genuine cycles at the expense of noise-induced ones.
The effect is data-dependent and non-monotone, as demonstrated by the divergent LUAD/GSE161711 results across the geometry family.
\end{remark}


% ============================================================
\section{Backend Implementation Details}
\label{sec:supp_backends}
% ============================================================

This section provides the full implementation specifications for each geometry backend, summarized in the main text (Section~2.4).

\subsection{Euclidean backend}

Standard $\ell_2$ distance in PCA space: $D^{(\text{eu})}_{ij} = \|x_i - x_j\|_2$.
Applied after optional PCA dimensionality reduction.
Advantages: simplicity, interpretability, computational efficiency.
Limitation: treats all directions equally and does not adapt to local density or curvature.

\subsection{Ollivier--Ricci curvature backend}

\begin{definition}[Ollivier--Ricci curvature on graphs \cite{ollivier2007,ollivier2009}]\label{def:supp_ricci}
Let $G = (V, E, w)$ be a weighted graph.
For each vertex $v$, define a probability measure $\mu_v$ supported on its neighbors (e.g., uniform on the $k$-nearest neighbors).
The \emph{Ollivier--Ricci curvature} of an edge $(u,v) \in E$ is
\[
\kappa(u,v) = 1 - \frac{W_1(\mu_u, \mu_v)}{d(u,v)},
\]
where $W_1$ denotes the Wasserstein-1 (earth mover's) distance.
\end{definition}

Implementation procedure:
\begin{enumerate}[leftmargin=1.5em]
\item Construct $k$-nearest-neighbor graph $G = (V, E, w)$ with $w_{ij} = \|x_i - x_j\|_2$.
\item For $T$ iterations: compute Ollivier--Ricci curvature $\kappa(i,j)$ for each edge via optimal transport between neighbor distributions; update weights $w_{ij}^{(t+1)} = w_{ij}^{(t)} - \eta \cdot \kappa^{(t)}(i,j) \cdot w_{ij}^{(t)}$.
\item Compute shortest-path distances on the deformed graph.
\end{enumerate}
Default parameters: $k=10$, $T=5$, $\eta=0.5$.

Positive curvature indicates convergent geometry (neighborhoods of $u$ and $v$ are close); negative curvature indicates divergent geometry characteristic of bridges, bottlenecks, or saddle-like structure between distinct communities.
Ricci flow upweights negatively curved (bridge) edges, delaying merger of distinct components in the filtration, preserving genuine cycles at the expense of noise-induced shortcuts.

\begin{remark}[Non-monotonicity of Ricci benefits]\label{rem:supp_ricci_nonmono}
Ricci flow does not uniformly improve downstream performance.
It helps when curvature correction suppresses misleading shortcut edges (as on LUAD, where Ricci achieves F1 $= 0.5012$ vs Euclidean $= 0.4342$).
It does not help when Euclidean neighborhoods already capture the relevant geometry (as on GSE161711, where Euclidean achieves F1 $= 0.7483$ vs Ricci $= 0.5944$).
This non-monotonicity motivates treating geometry as a diagnostic axis.
\end{remark}

\subsection{Diffusion backend}

\begin{definition}[Diffusion distance \cite{coifman2006diffusion}]\label{def:supp_diffusion}
Let $W$ be a weight matrix derived from a Gaussian kernel $W_{ij} = \exp(-\|x_i - x_j\|^2 / \sigma^2)$, and let $P = D_W^{-1} W$ be the row-normalized Markov transition matrix where $D_W = \operatorname{diag}(\sum_j W_{ij})$.
The \emph{diffusion distance} at time $t$ is
\[
D^{(\text{diff})}_{t}(i,j) = \Bigl( \sum_{\ell \ge 1} e^{-2t\mu_\ell} (\psi_\ell(i) - \psi_\ell(j))^2 \Bigr)^{1/2},
\]
where $\mu_\ell$ and $\psi_\ell$ are eigenvalues and right eigenvectors of $P$.
\end{definition}

The diffusion time parameter $t$ controls the scale of geometry: small $t$ recovers local Euclidean-like distances; large $t$ emphasizes global connectivity.
Diffusion distances can equivalently be expressed through the heat kernel $H_t = e^{-tL}$ as $D^{(\text{diff})}_t(i,j)^2 = H_t(i,i) + H_t(j,j) - 2H_t(i,j)$, connecting persistent homology to graph spectral methods.

Implementation: truncated eigendecomposition via \texttt{scipy.sparse.linalg.eigsh} with 50 eigenvectors.

\subsection{PHATE-like backend}

\begin{definition}[PHATE-like potential distance \cite{moon2019phate}]\label{def:supp_phate}
Let $P^t$ be the $t$-step transition matrix.
The \emph{potential representation} of sample $i$ is $U_i = -\log(P^t_{i,\cdot} + \epsilon)$, where $\epsilon > 0$ is a regularization constant.
The \emph{potential distance} is
\[
D^{(\text{phate})}(i,j) = \|U_i - U_j\|_2.
\]
\end{definition}

The log transformation converts transition probabilities into surprisal (information content).
The potential distance therefore measures how differently the random walks emanating from $i$ and $j$ explore the graph, related to the symmetrized KL divergence between diffusion distributions.
In information-geometric terms, $D^{(\text{phate})}$ approximates a Fisher--Rao-like distance on the statistical manifold of diffusion distributions.

Implementation: lightweight potential-distance computation without the full MDS step; $\epsilon = 10^{-7}$ for numerical stability.
Results may differ from those obtained with the complete PHATE pipeline including MDS dimensionality reduction.

\subsection{Distance-to-measure (DTM) backend}

\begin{definition}[Distance to measure \cite{chazal2011dtm, chazal2017robust}]\label{def:supp_dtm}
For a point cloud $X$ with empirical measure $\mu_X = \frac{1}{n}\sum_{i=1}^n \delta_{x_i}$ and a mass parameter $m_0 \in (0,1]$, the \emph{distance-to-measure function} at point $x$ is
\[
d_{\mu_X, m_0}(x) = \Bigl( \frac{1}{\lceil m_0 n \rceil} \sum_{j=1}^{\lceil m_0 n \rceil} \|x - x_{(j)}\|^2 \Bigr)^{1/2},
\]
where $x_{(1)}, \ldots, x_{(n)}$ are the points of $X$ sorted by distance to $x$.
\end{definition}

\begin{proposition}[Stability of DTM \cite{chazal2011dtm}]\label{prop:supp_dtm}
If $\mu$ and $\nu$ are two probability measures on $\mathbb{R}^p$, then
\[
\|d_{\mu, m_0} - d_{\nu, m_0}\|_\infty \le \frac{1}{\sqrt{m_0}} W_2(\mu, \nu),
\]
where $W_2$ is the Wasserstein-2 distance.
\end{proposition}

DTM-weighted distances delay inclusion of outlier-adjacent simplices, suppressing topological noise from density spikes.
Implementation: $k$NN-based measure estimation for computational efficiency.


% ============================================================
\section{Computational Timing Profiles}
\label{sec:supp_timing}
% ============================================================

\subsection{Phase 1: Pilot benchmark timing}

Table~\ref{tab:supp_timing} reports per-backend runtimes for the Phase~1 pilot benchmark (LUAD pilot $n=60$, GSE161711 $n=96$, GSE165087 bounded $n=320$) on an HPC cluster (48\,GB, 8 CPUs).

\begin{table}[h]
\centering
\small
\begin{tabular}{llr}
\toprule
Dataset & Backend & Seconds \\
\midrule
LUAD ($n=60$) & euclidean & 0.5 \\
LUAD & ricci & 5.2 \\
LUAD & diffusion & 0.5 \\
LUAD & phate\_like & 0.2 \\
LUAD & dtm & 0.2 \\
\midrule
GSE161711 ($n=96$) & euclidean & 8.0 \\
GSE161711 & ricci & 29.0 \\
GSE161711 & diffusion & 2.5 \\
GSE161711 & phate\_like & 1.2 \\
GSE161711 & dtm & 14.4 \\
\midrule
GSE165087 ($n=320$) & euclidean & 1099.0 \\
GSE165087 & ricci & 6.2 \\
GSE165087 & diffusion & 1.9 \\
GSE165087 & phate\_like & 36.3 \\
GSE165087 & dtm & 3932.1 \\
\bottomrule
\end{tabular}
\caption{Per-backend runtime in seconds for the Phase~1 pilot benchmark (HPC; 48\,GB, 8 CPUs).
These three datasets completed in approximately 3.3 hours total.
The DTM backend on GSE165087 is the most expensive computation at ${\sim}3932$ seconds due to $k$NN-based measure estimation on $n = 320$ cells.}
\label{tab:supp_timing}
\end{table}


% ============================================================
\section{Computational Scaling Analysis}
\label{sec:supp_scaling}
% ============================================================

The full-scale LUAD evaluation was conducted on a local workstation (Intel i9-9920X, 24 threads @ 3.50\,GHz, 125\,GB RAM).
Table~\ref{tab:supp_scaling} reports per-backend walltimes, revealing dramatic nonlinear scaling from pilot to full scale.

\begin{table}[h]
\centering\small
\begin{tabular}{lrrrl}
\toprule
Backend & Pilot ($n=60$) & Full ($n=275$) & Scaling & Bottleneck \\
\midrule
euclidean & 0.5\,s & 2,527\,s (42\,m) & ${\sim}5054\times$ & CVXPY cycle opt. \\
ricci & 5.2\,s & 8,896\,s (2.5\,h) & ${\sim}1711\times$ & Curvature flow + CVXPY \\
diffusion & 0.5\,s & 7,851\,s (2.2\,h) & ${\sim}15701\times$ & Spectral decomp. + CVXPY \\
phate\_like & 0.2\,s & 2,100\,s (35\,m) & ${\sim}10500\times$ & Log-potential + CVXPY \\
dtm & 0.2\,s & 6,487\,s (1.8\,h) & ${\sim}32435\times$ & $k$NN measure + CVXPY \\
\bottomrule
\end{tabular}
\caption{Walltime scaling from pilot ($n = 60$, 2-split, HPC 8 CPUs) to full scale ($n = 275$, 5-fold, local i9-9920X 24 threads).
The superlinear scaling reflects three compounding factors: (i) the distance matrix grows as $O(n^2)$; (ii) the Rips complex grows combinatorially, increasing both persistence computation and boundary-matrix size for CVXPY; (iii) 5-fold CV vs 2-split increases solver invocations by ${\sim}2.5\times$.
TIP gene extraction adds an additional 15--150 minutes per backend (not included).}
\label{tab:supp_scaling}
\end{table}

The extreme scaling factors (up to $32{,}000\times$ for DTM) underscore that CC-SPR's cycle optimization is the computational bottleneck, not the geometry backend itself.
The CVXPY elastic-net solver operates on boundary matrices whose size grows with the number of edges in the Rips complex, which scales superlinearly with $n$.
Future work on scalability should target approximate cycle representatives or subsampled boundary matrices rather than backend optimization alone.

\paragraph{Practical recommendation.}
For exploratory analyses, the pilot protocol ($n \leq 100$, 2--3 splits) completes in seconds and suffices for backend selection.
For definitive results, the full-scale protocol ($n = 275$, 5-fold CV) requires 2--3 hours per backend on a modern workstation and should be parallelized across backends.


% ============================================================
\section{Resource-Bounded Single-Cell Validation Details}
\label{sec:supp_scrna}
% ============================================================

Three larger GSE165087 configurations failed under cluster memory limits, establishing the resource boundary for single-cell analysis.
Table~\ref{tab:supp_scrna} reports the full execution envelope.

\begin{table}[h]
\centering
\small
\begin{tabular}{lllll}
\toprule
Run type & State & Elapsed & Requested memory & Peak RSS \\
\midrule
scRNA full & OUT\_OF\_MEMORY & 16:43:25 & 96\,GB & 100.3\,GB \\
scRNA fast-track & TIMEOUT & 20:00:06 & 96\,GB & 75.2\,GB \\
scRNA low-memory & OUT\_OF\_MEMORY & 11:25:59 & 48\,GB & 49.7\,GB \\
scRNA manuscript-safe & COMPLETED & 00:07:49 & 32\,GB & 2.9\,GB \\
scRNA ablation-safe & COMPLETED & --- & 32\,GB & ${\sim}$9.3\,GB \\
\bottomrule
\end{tabular}
\caption{Execution envelope for single-cell experiments on GSE165087.
The manuscript-safe configuration ($n=320$ cells) produces F1 $= 0.925$ for all methods (ceiling effect).
The ablation-safe configuration ($n=400$) shows CC-SPR ($0.823$) exceeding harmonic ($0.798$) while trailing standard ($0.888$).}
\label{tab:supp_scrna}
\end{table}


% ============================================================
\section{High-Bootstrap Stability and Arabidopsis Feasibility}
\label{sec:supp_feasibility}
% ============================================================

This section documents the higher-$n_{\text{boot}}$ stability experiments reported in the main text and explains the computational limits that prevented matching $n_{\text{boot}}$ across datasets.

\subsection{GSE161711 at $n_{\text{boot}} = 1000$}

All five backends were run on 5{,}000 variable genes with $n_{\text{boot}} = 1000$ and 0.8 subsample fraction.
Completion rates, concentration metrics, and per-backend wall times are reproduced in the main text (Table~\ref{tab:gse_boot1000} of the main manuscript).
The label-permutation test used the Euclidean backend, $n_{\text{perm}} = 200$, $n_{\text{splits}} = 5$, and $n_{\text{boot}} = 5$ per permutation.
Observed F1 $= 0.6876$; zero of 200 permutations matched or exceeded it, yielding an empirical $p$-value of $0.0000$ (95\% upper CI $\approx 0.015$ under the rule-of-three approximation).
Total wall time across the five backends plus permutation test was approximately 8.3 hours on a 12-core/24-thread Intel i9-9920X with 125~GB RAM (CPU only).

\subsection{Arabidopsis scalability limits}

The original Arabidopsis benchmark in the main text uses $n_{\text{boot}} = 3$ within 3-fold CV on 500 cells.
For the stability study we attempted to scale to $n_{\text{boot}} = 1000$ following the GSE161711 protocol.
Three successive attempts revealed a hard scalability frontier.

\begin{itemize}
\item \textbf{Attempt 1 ($n = 2{,}999$ cells, $n_{\text{boot}} = 20$)}: OOM-killed at $\sim 71$~GB peak RAM on the 125~GB workstation during the first timing-probe iteration. Root cause: the Rips complex on a 2{,}399-cell subsample in 50 PCA dimensions stores $O(n^3) \approx 1.4 \times 10^{10}$ potential triangles, dominating memory.

\item \textbf{Attempt 2 ($n = 1{,}000$ cells, $n_{\text{boot}} = 20$)}: timing probe completed with stable memory at $\sim 14$~GB, but per-iteration cost reached $12{,}084$~s ($3.4$~h). Extrapolating, $n_{\text{boot}} = 20$ would take $\sim 67$ hours. The CVXPY solver emitted ``Solution may be inaccurate'' warnings, indicating ill-conditioned boundary matrices on the dense PCA distance matrix. Killed after the timing probe.

\item \textbf{Attempt 3 ($n = 500$ cells, $n_{\text{boot}} = 50$)}: initial parallel runs on all five backends hung without producing checkpoints after more than 20 hours of accumulated CPU time; probe timings underestimated realized per-iteration cost by a factor of 20--30$\times$. Killed.

\item \textbf{Attempt 3b ($n = 500$ cells, $n_{\text{boot}} = 10$, Euclidean + DTM)}: \emph{succeeded}. Euclidean: 10/10 iterations, support = 46, entropy = 3.655, Gini = 0.323, elapsed 10{,}189~s ($2.83$~h). DTM: 10/10, support = 48, entropy = 3.733, Gini = 0.289, elapsed 11{,}609~s ($3.22$~h). Total wall time $\approx 6$~h.
\end{itemize}

\paragraph{Why Arabidopsis is pathologically slower than GSE161711.}
GSE161711 is a bulk dataset with $n = 96$ samples in a sparse gene-space distance matrix.
Arabidopsis is represented in 50 PCA dimensions, producing a \emph{dense} distance matrix where every cell pair is non-trivially close.
This has three consequences:
(i)~many edges and triangles enter the Rips filtration at low scales, inflating the boundary matrix passed to CVXPY;
(ii)~the elastic-net cycle problem is ill-conditioned on dense boundary matrices, triggering many solver iterations;
(iii)~per-iteration cost varies dramatically with the random subsample geometry.
The backend slowdown versus GSE161711 ranges from $\sim 9\times$ (DTM) to $\sim 227\times$ (Diffusion).

\paragraph{Why a reduced cell count is still a valid stability test.}
Reducing from 2{,}999 $\to$ 500 cells yields a \emph{conservative} stability assessment: fewer cells produce noisier geometry estimates, so TIP scores have higher bootstrap-to-bootstrap variance, making it harder (not easier) to demonstrate stability.
Any stability observed at 500 cells is therefore a lower bound on the method's stability at full scale.

\paragraph{Methodological implication.}
The Arabidopsis scalability frontier is fundamental to the exact-Rips + elastic-net approach, not an implementation bug.
Scaling CC-SPR to large single-cell atlases will require either approximate persistent homology (sparse Rips, witness complexes) or better-conditioned cycle optimizers (for example, L1 regression or first-order methods), both identified as future work.


% ============================================================
\section{Software Versions}
\label{sec:supp_software}
% ============================================================

\begin{table}[h]
\centering
\small
\begin{tabular}{ll}
\toprule
Package & Version \\
\midrule
Python & $\geq$ 3.10 \\
GUDHI & $\geq$ 3.8 \\
CVXPY & $\geq$ 1.4 \\
GraphRicciCurvature & $\geq$ 0.5 \\
scanpy & $\geq$ 1.9 \\
scikit-learn & $\geq$ 1.3 \\
gseapy & $\geq$ 1.1 \\
numpy & $\geq$ 1.24 \\
\bottomrule
\end{tabular}
\caption{Key software dependencies. A frozen environment file (\texttt{requirements-frozen.txt}) with exact versions is provided in the repository.}
\label{tab:supp_software}
\end{table}


% ============================================================
%  BIBLIOGRAPHY
% ============================================================
\bibliographystyle{unsrt}
\begin{thebibliography}{15}

\bibitem{harmonicPH}
D.~Gurnari, A.~Guzman-Saenz, F.~Utro, A.~Bose, S.~Basu, and L.~Parida.
Probing omics data via harmonic persistent homology.
\emph{Scientific Reports}, 15:3422, 2025.

\bibitem{friedman1998}
J.~Friedman.
Computing Betti numbers via combinatorial Laplacians.
\emph{Algorithmica}, 21(4):331--346, 1998.

\bibitem{zomorodian2005}
A.~Zomorodian and G.~Carlsson.
Computing persistent homology.
\emph{Discrete \& Computational Geometry}, 33(2):249--274, 2005.

\bibitem{zou2005elasticnet}
H.~Zou and T.~Hastie.
Regularization and variable selection via the elastic net.
\emph{JRSS Series B}, 67(2):301--320, 2005.

\bibitem{chazal2014persistence}
F.~Chazal, V.~de Silva, M.~Glisse, and S.~Oudot.
The structure and stability of persistence modules.
\emph{Springer}, 2016.

\bibitem{ollivier2007}
Y.~Ollivier.
Ricci curvature of metric spaces.
\emph{Comptes Rendus Mathematique}, 345(11):643--646, 2007.

\bibitem{ollivier2009}
Y.~Ollivier.
Ricci curvature of Markov chains on metric spaces.
\emph{Journal of Functional Analysis}, 256(3):810--864, 2009.

\bibitem{ni2019community}
C.-C.~Ni, Y.-Y.~Lin, F.~Luo, and J.~Gao.
Community detection on networks with Ricci flow.
\emph{Scientific Reports}, 9:9984, 2019.

\bibitem{coifman2006diffusion}
R.~R.~Coifman and S.~Lafon.
Diffusion maps.
\emph{Applied and Computational Harmonic Analysis}, 21(1):5--30, 2006.

\bibitem{moon2019phate}
K.~R.~Moon, D.~van Dijk, Z.~Wang, et al.
Visualizing structure and transitions in high-dimensional biological data.
\emph{Nature Biotechnology}, 37(12):1482--1492, 2019.

\bibitem{chazal2011dtm}
F.~Chazal, D.~Cohen-Steiner, and Q.~M\'erigot.
Geometric inference for probability measures.
\emph{Foundations of Computational Mathematics}, 11(6):733--751, 2011.

\bibitem{chazal2017robust}
F.~Chazal, B.~Fasy, F.~Lecci, B.~Michel, A.~Rinaldo, and L.~Wasserman.
Robust topological inference: distance to a measure and kernel distance.
\emph{JMLR}, 18(159):1--40, 2017.

\bibitem{haghverdi2016}
L.~Haghverdi, M.~B\"uttner, F.~A.~Wolf, F.~Buettner, and F.~J.~Theis.
Diffusion pseudotime robustly reconstructs lineage branching.
\emph{Nature Methods}, 13(10):845--848, 2016.

\end{thebibliography}

\end{document}
